\documentclass[11pt]{article}
\usepackage[margin=1in]{geometry}
\usepackage{graphicx}
\usepackage{tikz-cd}
\usepackage{multicol}
\usepackage{natbib}
\usepackage[spanish]{babel}
\usepackage[utf8]{inputenc}
\usepackage{amsmath}
\usepackage{amsfonts} 
\usepackage{amssymb}   
\setlength{\parindent}{0pt}
\setlength{\parskip}{1\baselineskip}

\title{\Large\textbf{Modelo pensional de Asofondos} \\[0.2cm]}
\author{\normalsize Colaboración Asofondos -- CARF -- Fedesarrollo}
\date{\normalsize \today}
\begin{document}
\maketitle

Este documento expone la arquitectura matemática y el diseño conceptual del modelo de microsimulación del sistema pensional colombiano, basado en \citet{becerra2025simulation}. El modelo está diseñado para superar las limitaciones de los análisis estáticos, al proyectar de manera intertemporal la evolución de los individuos, sus trayectorias laborales y contributivas, y las reglas institucionales que determinan el acceso y el monto de los beneficios pensionales.

A partir de esta estructura, el modelo permite evaluar la sostenibilidad financiera del sistema, estimando el déficit anual y el Valor Presente Neto (VPN) del pasivo pensional para el periodo 2025-2100, así como la evolución de los saldos del Fondo de Garantía de Pensión Mínima (FGPM) y del Fondo de Ahorro del Pilar Contributivo. Además, incorpora dimensiones de cobertura, suficiencia y equidad, al calcular tasas de reemplazo, pensiones promedio por pilar, brechas de género en densidades de cotización, subsidios implícitos por decil de ingreso y medidas distributivas como el coeficiente de Gini de pensiones.

El modelo también permite simular la coexistencia entre el marco de la Ley 100 de 1993 y el esquema de pilares introducido por la Ley 2381 de 2024, así como evaluar escenarios alternativos ante cambios demográficos, laborales, macroeconómicos o institucionales. Su implementación modular en Python lo convierte en un laboratorio prospectivo para analizar reformas paramétricas, choques macroeconómicos y la eventual incorporación de regímenes especiales como el FOMAG y CASUR.

Para materializar estas métricas agregadas a partir de microdatos, el marco conceptual se fundamenta en la evolución intertemporal de arreglos multidimensionales, o tensores, cuya estructura matemática se detalla a continuación.

\section{Estructura}

\subsection{Demografía y educación}
El punto de partida del modelo es un tensor poblacional:
\begin{align*}
\mathsf{N}
=
[n_{c,g,e,t}]
\in
\mathbb{R}_{\geq 0}^{|\mathcal{C}| \times |\mathcal{G}| \times |\mathcal{E}| \times |\mathcal{T}|}.
\end{align*}

Los índices del tensor se definen a partir de los siguientes conjuntos:
\begin{align*}
\mathcal{C} &= \{c_0,c_1,\ldots,c_\infty\}, \\
\mathcal{G} &= \{h,m\}, \\
\mathcal{E} &= \{e_L,e_M,e_H\}, \\
\mathcal{T} &= \{t_0,t_1,\ldots,t_\infty\}.
\end{align*}

Los componentes de estos conjuntos se interpretan como:
\begin{align*}
h &= \text{hombre}, 
& m &= \text{mujer}, \\[0.2cm]
e_L &= \text{logro educativo bajo}, 
& e_M &= \text{logro educativo medio}, 
& e_H &= \text{logro educativo alto}.
\end{align*}

Por tanto, cada elemento \(n_{c,g,e,t}\) representa el número de personas vivas pertenecientes a la cohorte \(c\), de sexo \(g\), con logro educativo \(e\), en el año \(t\). Asimismo, para un sexo fijo \(\hat{g}\in\mathcal{G}\) y un año fijo \(\hat{t}\in\mathcal{T}\), la rebanada del tensor correspondiente a \((\hat{g},\hat{t})\) está dada por:
\[
\mathbf{N}_{\hat{g},\hat{t}}
=
[n_{c,\hat{g},e,\hat{t}}]
\in
\mathbb{R}_{\geq 0}^{|\mathcal{C}| \times |\mathcal{E}|}.
\]

Esta rebanada es equivalente a una matriz de cohortes por logro educativo:
\begin{align*}
\mathbf{N}_{\hat{g},\hat{t}}
=
\begin{pmatrix}
n_{c_0,\hat{g},e_L,\hat{t}} 
& n_{c_0,\hat{g},e_M,\hat{t}} 
& n_{c_0,\hat{g},e_H,\hat{t}} \\[0.2cm]
n_{c_1,\hat{g},e_L,\hat{t}} 
& n_{c_1,\hat{g},e_M,\hat{t}} 
& n_{c_1,\hat{g},e_H,\hat{t}} \\
\vdots & \vdots & \vdots \\
n_{c_\infty,\hat{g},e_L,\hat{t}} 
& n_{c_\infty,\hat{g},e_M,\hat{t}} 
& n_{c_\infty,\hat{g},e_H,\hat{t}}
\end{pmatrix}.
\end{align*}




\subsection{Mercado laboral}

Para incorporar la dinámica económica que antecede a los resultados pensionales de los individuos, así como capturar las particularidades de los regímenes especiales y modelar los riesgos de invalidez, la población contenida en cada celda del tensor \(\mathsf{N}\) se descompone según su condición en el mercado laboral. Se define el conjunto ampliado de estados laborales posibles como:
\begin{align*}
\mathcal{L} &= \{a, i, d, n, x, p, v\},
\end{align*}
donde los componentes se interpretan de la siguiente manera:
\begin{align*}
a &= \text{asalariado (cotizante dependiente)}, \\
i &= \text{independiente (cuenta propia, cotizante o informal)}, \\
d &= \text{desempleado}, \\
n &= \text{inactivo (fuera de la fuerza laboral, sin pensión de invalidez)}, \\
x &= \text{inactivo discapacitado (elegible para o recibiendo pensión de invalidez)}, \\
p &= \text{maestro (régimen especial del Magisterio)}, \\
v &= \text{militar (régimen especial de la Fuerza Pública)}.
\end{align*}

Para cada celda poblacional \((c,g,e,t)\), se define el vector fila de distribución laboral:
\begin{align*}
\boldsymbol{\pi}_{c,g,e,t}
=
\begin{pmatrix}
\pi_{c,g,e,a,t} &
\pi_{c,g,e,i,t} &
\pi_{c,g,e,d,t} &
\pi_{c,g,e,n,t} &
\pi_{c,g,e,x,t} &
\pi_{c,g,e,p,t} &
\pi_{c,g,e,v,t}
\end{pmatrix}
\in [0,1]^{1\times 7},
\end{align*}
donde
\begin{align*}
\sum_{l\in\mathcal{L}} \pi_{c,g,e,l,t}=1
\qquad
\forall (c,g,e,t).
\end{align*}

La población en cada estado laboral se obtiene al multiplicar la población total de la celda por su distribución laboral:
\begin{align*}
m_{c,g,e,l,t}
=
n_{c,g,e,t}\,\pi_{c,g,e,l,t}
\qquad
\forall l\in\mathcal{L}.
\end{align*}
En forma vectorial:
\begin{align*}
\mathbf{m}_{c,g,e,t}
=
n_{c,g,e,t}\boldsymbol{\pi}_{c,g,e,t}
\in
\mathbb{R}_{\geq 0}^{1\times 7}.
\end{align*}

De este modo, el tensor general de la población se expande a un tensor de orden cinco:
\begin{align*}
\mathsf{M}
=
[m_{c,g,e,l,t}]
\in
\mathbb{R}_{\geq 0}^{|\mathcal{C}| \times |\mathcal{G}| \times |\mathcal{E}| \times |\mathcal{L}| \times |\mathcal{T}|}.
\end{align*}

Por construcción, este tensor conserva la consistencia agregada con respecto al tensor poblacional original \(\mathsf{N}\):
\begin{align*}
n_{c,g,e,t}
=
\sum_{l \in \mathcal{L}} m_{c,g,e,l,t}
\qquad
\forall (c,g,e,t)
\in
\mathcal{C}\times\mathcal{G}\times\mathcal{E}\times\mathcal{T}.
\end{align*}

La transición intertemporal de la distribución laboral entre el año \(t\) y el año \(t+1\) se modela como un proceso de Markov finito, gobernado por una matriz de transición estocástica por filas. Para cada cohorte \(c\), sexo \(g\), logro educativo \(e\) y año \(t\), se define:
\begin{align*}
\mathbf{P}_{c,g,e,t}
=
\left[p_{jk}(c,g,e,t)\right]_{j,k\in\mathcal{L}}
\in
[0,1]^{7\times 7}.
\end{align*}

Dado el ordenamiento \(\mathcal{L}=\{a,i,d,n,x,p,v\}\), esta matriz puede escribirse como:
\begin{align*}
\mathbf{P}_{c,g,e,t}
=
\begin{pmatrix}
p_{aa} & p_{ai} & p_{ad} & p_{an} & p_{ax} & p_{ap} & p_{av} \\
p_{ia} & p_{ii} & p_{id} & p_{in} & p_{ix} & p_{ip} & p_{iv} \\
p_{da} & p_{di} & p_{dd} & p_{dn} & p_{dx} & p_{dp} & p_{dv} \\
p_{na} & p_{ni} & p_{nd} & p_{nn} & p_{nx} & p_{np} & p_{nv} \\
p_{xa} & p_{xi} & p_{xd} & p_{xn} & p_{xx} & p_{xp} & p_{xv} \\
p_{pa} & p_{pi} & p_{pd} & p_{pn} & p_{px} & p_{pp} & p_{pv} \\
p_{va} & p_{vi} & p_{vd} & p_{vn} & p_{vx} & p_{vp} & p_{vv}
\end{pmatrix}_{c,g,e,t}.
\end{align*}

Cada elemento \(p_{jk}(c,g,e,t)\) denota la probabilidad de que un individuo de la cohorte \(c\), sexo \(g\) y logro educativo \(e\) transite del estado laboral \(j\in\mathcal{L}\) al estado laboral \(k\in\mathcal{L}\) entre los años \(t\) y \(t+1\). Por definición, las filas satisfacen:
\begin{align*}
\sum_{k \in \mathcal{L}} p_{jk}(c,g,e,t)=1
\qquad
\forall j\in\mathcal{L}.
\end{align*}

La evolución de la distribución laboral dentro de cada celda poblacional está dada por:
\begin{align*}
\boldsymbol{\pi}_{c,g,e,t+1}
=
\boldsymbol{\pi}_{c,g,e,t}
\mathbf{P}_{c,g,e,t}.
\end{align*}

Finalmente, una vez determinada la distribución laboral del año \(t+1\), los conteos por estado laboral se obtienen usando la población proyectada de la celda:
\begin{align*}
\mathbf{m}_{c,g,e,t+1}
=
n_{c,g,e,t+1}
\boldsymbol{\pi}_{c,g,e,t+1}.
\end{align*}

\subsection{Sistema pensional}

\bibliographystyle{apalike}
\bibliography{Referencias}

\end{document}